\chapter{Практическая часть}

\section{Исходные данные конкретного варианта}

Матрица стоимостей варианта \textbf{№8}:
$$
    C = 
    \begin{bmatrix}
        7 & 4 & 3 & 8 & 2 \\
        4 & 5 & 1 & 6 & 3 \\
        8 & 4 & 5 & 7 & 2 \\
        1 & 2 & 4 & 7 & 2 \\
        3 & 9 & 9 & 2 & 5
    \end{bmatrix}. 
$$

\section{Текст программы}

Текст программы представлен в листинге~\ref{lst:text}. 

\mylisting[./code/lst.m]{Текст программы}{}{label=lst:text}

\section{Результаты работы программы}

\begin{center}
\captionsetup{justification=raggedright,singlelinecheck=off}
\begin{lstlisting}[label=lst:lab]
Матрица стоимостей:
     7     4     3     8     2
     4     5     1     6     3
     8     4     5     7     2
     1     2     4     7     2
     3     9     9     2     5

Подготовительный этап:

Минимальные элементы столбцов:
     1     2     1     2     2

Матрица стоимостей после вычитания минимальных элементов из столбцов:
     6     2     2     6     0
     3     3     0     4     1
     7     2     4     5     0
     0     0     3     5     0
     2     7     8     0     3

Минимальные элементы строк:
     0
     0
     0
     0
     0

Матрица стоимостей после вычитания минимальных элементов из строк:
     6     2     2     6     0
     3     3     0     4     1
     7     2     4     5     0
     0     0     3     5     0
     2     7     8     0     3

Начальная СНН:
	6	2	2	6	0*
	3	3	0*	4	1
	7	2	4	5	0
	0*	0	3	5	0
	2	7	8	0*	3
Колическтво нулей СНН: k = 4

Основной этап:

k = 4 < n = 5. Требуется 1 итераций

Итерация №1

СНН с начально выделенными столбцами:
	6	2	2	6	0*
	3	3	0*	4	1
	7	2	4	5	0
	0*	0	3	5	0
	2	7	8	0*	3
	+	 	+	+	+

Отметим 0' найденный невыделенный 0:
	6	2	2	6	0*
	3	3	0*	4	1
	7	2	4	5	0
	0*	0'	3	5	0
	2	7	8	0*	3
	+	 	+	+	+

Перенесем выделение со столбца на строку, т.к. на строке с 0' есть 0*:
	6	2	2	6	0*
	3	3	0*	4	1
	7	2	4	5	0
	0*	0'	3	5	0	+
	2	7	8	0*	3
	 	 	+	+	+

Среди невыделенных элементов нет 0, преобразуем матрицу:
	4	0	2	6	0*
	1	1	0*	4	1
	5	0	4	5	0
	0*	0'	5	7	2	+
	0	5	8	0*	3
	 	 	+	+	+

Отметим 0' найденный невыделенный 0:
	4	0	2	6	0*
	1	1	0*	4	1
	5	0	4	5	0
	0*	0'	5	7	2	+
	0'	5	8	0*	3
	 	 	+	+	+

Перенесем выделение со столбца на строку, т.к. на строке с 0' есть 0*:
	4	0	2	6	0*
	1	1	0*	4	1
	5	0	4	5	0
	0*	0'	5	7	2	+
	0'	5	8	0*	3	+
	 	 	+	 	+

Отметим 0' найденный невыделенный 0:
	4	0'	2	6	0*
	1	1	0*	4	1
	5	0	4	5	0
	0*	0'	5	7	2	+
	0'	5	8	0*	3	+
	 	 	+	 	+

Перенесем выделение со столбца на строку, т.к. на строке с 0' есть 0*:
	4	0'	2	6	0*	+
	1	1	0*	4	1
	5	0	4	5	0
	0*	0'	5	7	2	+
	0'	5	8	0*	3	+
	 	 	+	 	 

Отметим 0' найденный невыделенный 0:
	4	0'	2	6	0*	+
	1	1	0*	4	1
	5	0'	4	5	0
	0*	0'	5	7	2	+
	0'	5	8	0*	3	+
	 	 	+	 	 

Построена L-цепочка: 0' (3, 2) 

X_opt = 
     0     0     0     0     1
     0     0     1     0     0
     0     1     0     0     0
     1     0     0     0     0
     0     0     0     1     0

f_opt = 10
\end{lstlisting}
\end{center}